\documentclass{article}

% these two are probably only necessary for the draft
\usepackage{geometry}
\usepackage{pdflscape}

% these are needed for substantive stuff
\usepackage{amsfonts}
\usepackage{amsmath}
\usepackage{amssymb}
\usepackage{amsthm}
\usepackage{stmaryrd}
\usepackage{ebproof}
\usepackage{mathtools}
\usepackage{simplebnf}

%%%%%%%%%%%%%%%%%%%%%%%%%
% GUIDE TO LATEX MACROS %
%%%%%%%%%%%%%%%%%%%%%%%%%

% There are two languages:
%   + 𝐁, the base language
%   + 𝐋, the logic language
% - \L<X> is the symbol for language X

% - There are five "semantics symbols":
%   + 𝔼, for exact values and evaluation semantics
%   + 𝔸, for approximation values
%   + ℂ, for clairvoyance semantics
%   + 𝔻, for demand semantics
%   + 𝕃, for logic semantics
% - \S<X> is the semantics symbol X

% - \T<X> is the type X
% - \S<X> is the set of values for the type X
% - \t<X> is the term (or term-level construct) X
% - \s<X> is the semantic construct <X>

% - \T<X> is also the typing relation for language X
% - \E<X> is the evaluation relation for semantics X
% - \trans<X> is the translation from B to L corresponding to semantics X

% - Macros are generally n-ary; i.e., the macro \TThunk for thunk types takes
% one argument, the type argument.  If we need a zero-arity version of macro M,
% we use <M>Z.  For example, the zero-ary version of \TThunk is \TThunkZ.

%%%%%%%%%%%%%%%%%%%%%
% MACRO DEFINITIONS %
%%%%%%%%%%%%%%%%%%%%%

% paired delimiters
\DeclarePairedDelimiter{\bbrackets}{\llbracket}{\rrbracket}
\DeclarePairedDelimiter{\braces}{\lbrace}{\rbrace}
\DeclarePairedDelimiter{\ceil}{\lceil}{\rceil}
\DeclarePairedDelimiter{\floor}{\lfloor}{\rfloor}
\DeclarePairedDelimiter{\verts}{\lvert}{\rvert}

% theorem environments
\newtheorem{lemma}{Lemma}
\newtheorem{theorem}{Theorem}

% math notation
\newcommand{\N}{\mathbb{N}}
\newcommand{\set}[1]{\braces*{#1}}
\newcommand{\setbuilder}[2]{\set{{#1} \mid {#2}}}

% language symbols
\newcommand{\LB}{\mathbf{B}}
\newcommand{\LL}{\mathbf{L}}

% semantics symbols
\newcommand{\SA}{\mathbb{A}}
\newcommand{\SC}{\mathbb{C}}
\newcommand{\SD}{\mathbb{D}}
\newcommand{\SE}{\mathbb{E}}
\newcommand{\SL}{\mathbb{L}}

% semantic mappings on types
\newcommand{\Approx}[1]{\SA\bbrackets{#1}}
\newcommand{\Eval}[1]{\SE\bbrackets{#1}}

% types
\newcommand{\TBool}{\texttt{Bool}}
\newcommand{\TProd}[2]{\texttt{\(#1\)\ *\ \(#2\)}}
\newcommand{\TNat}{\texttt{Nat}}
\newcommand{\TList}[1]{\texttt{List \({#1}\)}}
\newcommand{\TThunkZ}{\texttt{Thunk}}
\newcommand{\TThunk}[1]{\texttt{\TThunkZ\ \(#1\)}}
\newcommand{\TM}[1]{\texttt{M\ \(#1\)}}

% terms
\newcommand{\tlet}[3]{\texttt{let\ \({#1}\)\ =\ \({#2}\)\ in\ \({#3}\)}}
\newcommand{\ttrue}{\texttt{true}}
\newcommand{\tfalse}{\texttt{false}}
\newcommand{\teq}[2]{\texttt{\(#1\)\ ==\ \(#2\)}}
\newcommand{\tle}[2]{\texttt{\(#1\)\ <=\ \(#2\)}}
\newcommand{\tif}[3]{\texttt{if\ \({#1}\)\ then\ \({#2}\)\ else\ \({#3}\)}}
\newcommand{\tpair}[2]{\texttt{(\(#1\),\ \(#2\))}}
\newcommand{\tadd}[2]{\texttt{\(#1\)\ +\ \(#2\)}}
\newcommand{\tnil}{\texttt{[]}}
\newcommand{\tcons}[2]{\texttt{\(#1\)\ ::\ \(#2\)}}
\newcommand{\tfoldrZ}{\texttt{foldr}}
\newcommand{\tfoldr}[5]{\tfoldrZ\texttt{(\(\lambda{#1}{#2}\).\(#3\),\ \(#4\),\ \(#5\))}}
\newcommand{\tthunk}[1]{\texttt{thunk\ \({#1}\)}}
\newcommand{\tundefined}{\texttt{undefined}}
\newcommand{\tlazyZ}{\texttt{lazy}}
\newcommand{\tlazy}[1]{\texttt{\tlazyZ\ \(#1\)}}
\newcommand{\tforce}[1]{\texttt{force\ \({#1}\)}}
\newcommand{\tcase}[4]{\texttt{case \(#1\) of thunk \(#2\) => \(#3\); undefined => \(#4\)}}
\newcommand{\ttick}[1]{\texttt{tick\ \({#1}\)}}
\newcommand{\tchoose}[2]{\texttt{\(#1\)\ ?\ \(#2\)}}
\newcommand{\tfree}[1]{\texttt{free\ \(#1\)}}
\newcommand{\tfail}{\texttt{fail}}
\newcommand{\tassert}[1]{\texttt{assert\ \(#1\)}}
\newcommand{\treturn}[1]{\texttt{return\ \(#1\)}}
\newcommand{\ttranspose}[1]{\texttt{transpose\ \(#1\)}}
\newcommand{\tfoldMZ}{\texttt{foldM}}
\newcommand{\tfoldM}[5]{\texttt{\tfoldMZ(\(\lambda {#1}{#2}.{#3}\),\ \(#4\),\ \(#5\))}}

% semantic sets
\newcommand{\SBool}{\set{\strue, \sfalse}}

% semantic values and operations
\newcommand{\strue}{\text{true}}
\newcommand{\sfalse}{\text{false}}
\newcommand{\sle}[2]{{#1} \leqslant {#2}}
\newcommand{\snle}[2]{{#1} \nleqslant {#2}}
\newcommand{\spair}[2]{({#1}, {#2})}
\newcommand{\snil}{\varepsilon}
\newcommand{\scons}[2]{(#1, #2)}
\newcommand{\sfoldr}[3]{\text{foldr}({#1}, {#2}, {#3})}
\newcommand{\sthunk}[1]{\text{thunk\ \({#1}\)}}
\newcommand{\sundefined}{\bot}

% typing rules

% base
\newcommand{\TB}[3]{{#1} \vdash^\LB {#2} : {#3}}
% logic
\newcommand{\TL}[3]{{#1} \vdash^\LL {#2} : {#3}}

% semantics

% evaluation
\newcommand{\EE}[2]{\SE\bbrackets{#1}({#2})}
% clairvoyance
\newcommand{\EC}[4]{{{#1} \mathrel{;} {#2}} \Downarrow^\SC_{#3} {#4}}
% demand
\newcommand{\ED}[3]{\SD\bbrackets{#1}({#2}, {#3})}
% clairvoyance foldr
\newcommand{\ECF}[8]{{{#1} \mathrel{;} \lambda{#2}{#3}.{#4}, {#5}, {#6}} \Downarrow^{\text{\(\SC\)-\tfoldrZ}}_{#7} {#8}}
% logic
\newcommand{\EL}[3]{{{#1} \mathrel{;} {#2}} \Downarrow^\SL {#3}}
% logic foldr
\newcommand{\ELF}[7]{{{#1} \mathrel{;} \lambda{#2}{#3}.{#4}, {#5}, {#6}} \Downarrow^{\text{\(\SL\)-\tfoldrZ}} {#7}}
% logic foldM
\newcommand{\ELFM}[7]{{{#1} \mathrel{;} \lambda{#2}{#3}.{#4}, {#5}, {#6}} \Downarrow^{\text{\(\SL\)-\tfoldMZ}} {#7}}

% translations
\newcommand{\transC}[1]{\SC\floor{#1}}
\newcommand{\transD}[2]{\SD\floor{#1}({#2})}

% rules font
\newcommand{\rulename}[1]{\textsc{#1}}

\begin{document}

\section{The base language}

\subsection{Syntax}

\begin{center}
\begin{bnf}
\( A \) ::=
| \TBool
| \TList{A}
| \TThunk{A}
;;
\end{bnf}
\end{center}

\begin{center}
\begin{bnf}
$t$ ::=
| $x$
| \tlet{x}{t}{t}
| \ttrue
| \tfalse
| \tif{t}{t}{t}
| \tnil
| \tcons{t}{t}
| \tfoldr{x}{x}{t}{t}{t}
| \tlazy{t}
| \tforce{t}
| \ttick{t}
;;
\end{bnf}
\end{center}

Note that the language does not include functions.  This is due to a technical limitation of demand semantics.

\subsection{Typing}

\begin{center}
\begin{prooftree}
  \hypo{\Gamma(x) = A}
  \infer1{\Gamma \vdash^\LB x : A}
\end{prooftree}
\end{center}

\begin{center}
\begin{prooftree}
  \hypo{\Gamma \vdash^\LB t_1 : A}
  \hypo{\Gamma, x : A \vdash^\LB t_2 : B}
  \infer2{\Gamma \vdash^\LB \texttt{let $x$ = $t_1$ in $t_2$} : B}
\end{prooftree}
\end{center}

\begin{center}
\begin{prooftree}
  \infer0{\Gamma \vdash^\LB \ttrue : \TBool}
\end{prooftree}
\end{center}

\begin{center}
\begin{prooftree}
  \infer0{\Gamma \vdash^\LB \tfalse : \TBool}
\end{prooftree}
\end{center}

\begin{center}
\begin{prooftree}
  \hypo{\Gamma \vdash^\LB t_1 : \TBool}
  \hypo{\Gamma \vdash^\LB t_2 : A}
  \hypo{\Gamma \vdash^\LB t_3 : A}
  \infer3{\Gamma \vdash^\LB \tif{t_1}{t_2}{t_3} : A}
\end{prooftree}
\end{center}

\begin{center}
\begin{prooftree}
  \infer0{\Gamma \vdash^\LB \tnil : \TList{A}}
\end{prooftree}
\end{center}

\begin{center}
\begin{prooftree}
  \hypo{\Gamma \vdash^\LB t_1 : A}
  \hypo{\Gamma \vdash^\LB t_2 : \TThunk{(\TList{A})}}
  \infer2{\Gamma \vdash^\LB \tcons{t_1}{t_2} : \TList{A}}
\end{prooftree}
\end{center}

\begin{center}
\begin{prooftree}
  \hypo{\Gamma, x_1 : A, x_2 : \TThunk{B} \vdash^\LB t_1 : B}
  \hypo{\Gamma \vdash^\LB t_2 : B}
  \hypo{\Gamma \vdash^\LB t_3 : \TList{A}}
  \infer3{\Gamma \vdash^\LB \tfoldr{x_1}{x_2}{t_1}{t_2}{t_3} : B}
\end{prooftree}
\end{center}

\begin{center}
\begin{prooftree}
  \hypo{\Gamma \vdash^\LB t : A}
  \infer1{\Gamma \vdash^\LB \tlazy{t} : \TThunk{A}}
\end{prooftree}
\end{center}

\begin{center}
\begin{prooftree}
  \hypo{\Gamma \vdash^\LB t : \TThunk{A}}
  \infer1{\Gamma \vdash^\LB \tforce{t} : A}
\end{prooftree}
\end{center}

\begin{center}
\begin{prooftree}
  \hypo{\Gamma \vdash^\LB t : A}
  \infer1{\Gamma \vdash^\LB \ttick{t} : A}
\end{prooftree}
\end{center}

\subsection{Exact values and approximations}

\begin{center}
  \begin{align*}
    \Eval{\TBool} &= \SBool \\
    \Eval{\TThunk{A}} &= \Eval{A} \\
    \Eval{\TList{A}} &= \braces{\snil} \cup \setbuilder{\scons{v_1}{v_2}}{v_1 \in \Eval{A}, v_2 \in \Eval{\TList{A}}}
  \end{align*}

  \begin{align*}
    \Approx{\TBool} &= \SBool \\
    \Approx{\TThunk{A}} &= \set{\sundefined} \cup \setbuilder{\sthunk{v}}{v \in \Approx{A}} \\
    \Approx{\TList{A}} &= \set{\snil} \cup \setbuilder{\scons{v_1}{v_2}}{v_1 \in \Approx{A}, v_2 \in \Approx{\TThunk{(\TList{A})}}}
  \end{align*}
\end{center}

\subsection{Evaluation semantics}

\begin{align*}
  \EE{x}{\gamma} &= \gamma(x) \\
  \EE{\tlet{x}{t_1}{t_2}}{\gamma} &= \EE{t_2}{\gamma, x \mapsto \EE{t_1}{\gamma}} \\
  \EE{\ttrue}{\gamma} &= \strue \\
  \EE{\tfalse}{\gamma} &= \sfalse \\
  \EE{\tif{t_1}{t_2}{t_3}}{\gamma} &=
    \begin{cases}
      \EE{t_2}{\gamma} & \text{if \( \EE{t_{1}}{\gamma} = \strue \)} \\
      \EE{t_3}{\gamma} & \text{otherwise}
    \end{cases} \\
  \EE{\tnil}{\gamma} &= \varepsilon \\
  \EE{\tcons{t_1}{t_2}}{\gamma} &= \scons{\EE{t_1}{\gamma}}{\EE{t_2}{\gamma}} \\
  \EE{\tfoldr{x_1}{x_2}{t_1}{t_2}{t_3}}{\gamma} &= \sfoldr{(v_1, v_2) \mapsto \EE{t_1}{\gamma, x_1 \mapsto v_1, x_2 \mapsto v_2}}{\EE{t_2}{\gamma}}{\EE{t_3}{\gamma}} \\
  \EE{\tlazy{t}}{\gamma} &= \EE{t}{\gamma} \\
  \EE{\tforce{t}}{\gamma} &= \EE{t}{\gamma} \\
  \EE{\ttick{t}}{\gamma} &= \EE{t}{\gamma}
\end{align*}

\begin{align*}
  \sfoldr{f}{y}{(x_1, x_2, \ldots, x_n)} = f(x_1, f(x_2, \cdots f(x_n, y) \cdots))
\end{align*}

\begin{lemma}
  If \( \TB{\Gamma}{t}{A} \) and \( \gamma \in \Eval{\Gamma} \), then \( \EE{t}{\gamma} \in \Eval{A} \).
\end{lemma}

\subsection{Clairvoyance semantics}

Clairvoyance semantics has the judgment schema \( \EC{\gamma}{t}{c}{v} \).  Intuitively, \( c \) is the number of \( \texttt{tick} \)s encountered while evaluating \( t \) to \( v \).

\begin{center}
  \begin{prooftree}
    \hypo{\gamma(x) = v}
    \infer1{\EC{\gamma}{x}{0}{v}}
  \end{prooftree}
\end{center}

\begin{center}
  \begin{prooftree}
    \hypo{\EC{\gamma}{t_1}{c_1}{v_1}}
    \hypo{\EC{\gamma, x \mapsto v_1}{t_2}{c_2}{v_2}}
    \infer2{\EC{\gamma}{\tlet{x}{t_1}{t_2}}{c_1 + c_2}{v_2}}
  \end{prooftree}
\end{center}

\begin{center}
  \begin{prooftree}
    \infer0{\EC{\gamma}{\ttrue}{0}{\strue}}
  \end{prooftree}
\end{center}

\begin{center}
  \begin{prooftree}
    \infer0{\EC{\gamma}{\tfalse}{0}{\sfalse}}
  \end{prooftree}
\end{center}

\begin{center}
  \begin{prooftree}
    \hypo{\EC{\gamma}{t_1}{c_1}{\strue}}
    \hypo{\EC{\gamma}{t_2}{c_2}{v}}
    \infer2{\EC{\gamma}{\tif{t_1}{t_2}{t_3}}{c_1 + c_2}{v}}
  \end{prooftree}
\end{center}

\begin{center}
  \begin{prooftree}
    \hypo{\EC{\gamma}{t_1}{c_1}{\sfalse}}
    \hypo{\EC{\gamma}{t_3}{c_3}{v}}
    \infer2{\EC{\gamma}{\tif{t_1}{t_2}{t_3}}{c_1 + c_3}{v}}
  \end{prooftree}
\end{center}

\begin{center}
  \begin{prooftree}
    \infer0{\EC{\gamma}{\tnil}{0}{\snil}}
  \end{prooftree}
\end{center}

\begin{center}
  \begin{prooftree}
    \hypo{\EC{\gamma}{t_1}{c_1}{v_1}}
    \hypo{\EC{\gamma}{t_2}{c_1}{v_2}}
    \infer2{\EC{\gamma}{\tcons{t_1}{t_2}}{c_1 + c_2}{\scons{v_1}{v_2}}}
  \end{prooftree}
\end{center}

\begin{center}
  \begin{prooftree}
    \hypo{\EC{\gamma}{t_3}{c_1}{v_1}}
    \hypo{\ECF{\gamma}{x_1}{x_2}{t_1}{t_2}{v_1}{c_2}{v_2}}
    \infer2[\rulename{L-Foldr}]{\EC{\gamma}{\tfoldr{x_1}{x_2}{t_1}{t_2}{t_3}}{c_1 + c_2}{v_2}}
  \end{prooftree}
\end{center}

\begin{center}
  \begin{prooftree}
    \infer0{\EC{\gamma}{\tlazy{t}}{0}{\sundefined}}
  \end{prooftree}
\end{center}

\begin{center}
  \begin{prooftree}
    \hypo{\EC{\gamma}{t}{c}{v}}
    \infer1{\EC{\gamma}{\tlazy{t}}{c}{\sthunk{v}}}
  \end{prooftree}
\end{center}

\begin{center}
  \begin{prooftree}
    \hypo{\EC{\gamma}{t}{c}{\sthunk{v}}}
    \infer1{\EC{\gamma}{\tforce{t}}{c}{v}}
  \end{prooftree}
\end{center}

\begin{center}
  \begin{prooftree}
    \hypo{\EC{\gamma}{t}{c}{v}}
    \infer1{\EC{\gamma}{\ttick{t}}{c + 1}{v}}
  \end{prooftree}
\end{center}

Note that \( \tfoldrZ \) is evaluated by first evaluting the list argument to a value and then iterating over that value.  The iteration is expressed in the following auxiliary relation.
\begin{center}
  \begin{prooftree}
    \hypo{\EC{\gamma}{t_2}{c}{v}}
    \infer1[\rulename{C-Foldr-Nil}]{\ECF{\gamma}{x_1}{x_2}{t_1}{t_2}{\snil}{c}{v}}
  \end{prooftree}
\end{center}
\begin{center}
  \begin{prooftree}
    \hypo{\EC{\gamma, x_1 \mapsto v_1, x_2 \mapsto \sundefined}{t_1}{c}{v}}
    \infer1[\rulename{C-Foldr-Undefined}]{\ECF{\gamma}{x_1}{x_2}{t_1}{t_2}{\scons{v_1}{\sundefined}}{c}{v}}
  \end{prooftree}
\end{center}
\begin{center}
  \begin{prooftree}
    \hypo{\ECF{\gamma}{x_1}{x_2}{t_1}{t_2}{v_2}{c_1}{v_3}}
    \hypo{\EC{\gamma, x_1 \mapsto v_1, x_2 \mapsto \sthunk{v_3}}{t_1}{c_2}{v}}
    \infer2[\rulename{C-Foldr-Thunk}]{\ECF{\gamma}{x_1}{x_2}{t_1}{t_2}{\scons{v_1}{\sthunk{v_2}}}{c_1 + c_2}{v}}
  \end{prooftree}
\end{center}

\begin{lemma}
  If \( \TB{\Gamma}{t}{A} \), \( \gamma \in \Approx{\Gamma} \), and \( \EC{\gamma}{t}{\relax}{v} \), then \( v \in \Approx{A} \).
\end{lemma}

\begin{lemma}
  If \( \TB{\Gamma}{t}{A} \) and \( \gamma \in \Approx{\Gamma} \), then \( \EC{\gamma}{t}{\relax}{\relax} \).
\end{lemma}

\subsection{Demand semantics}

\begin{align*}
  \ED{x}{\gamma}{a} &= (\braces{x \mapsto a}, 0) \\
  \ED{\ttrue}{\gamma}{a} &= (\bot, 0) \\
  \ED{\tfalse}{\gamma}{a} &= (\bot, 0) \\
  \ED{\tnil}{\gamma}{a} &= (\bot, 0) \\
  \ED{\tcons{t_1}{t_2}}{\gamma}{a} &= \ED{t_1}{\gamma}{a} \sqcup \ED{t_2}{\gamma}{a}
\end{align*}

\section{The logic language}

\subsection{Syntax}

\begin{center}
\begin{bnf}
\( A \) ::=
| \TBool
| \TProd{A}{A}
| \TThunk{A}
| \TNat
| \TList{A}
;;
\end{bnf}
\end{center}

\begin{center}
\begin{bnf}
$t$ ::=
| $x$
| \tlet{x}{t}{t}
| \ttrue
| \tfalse
| \teq{t}{t}
| \tle{t}{t}
| \tif{t}{t}{t}
| \tpair{t}{t}
| \tthunk{t}
| \tundefined
| \tcase{t}{x}{t}{t}
| $n$
| \tadd{t}{t}
| \tif{t}{t}{t}
| \tnil
| \tcons{t}{t}
| \tfoldr{x}{x}{t}{t}{t}
| \tchoose{t}{t}
| \tfree{A}
| \tfail
;;
\end{bnf}
\end{center}

\subsection{Typing}

\begin{center}
  \begin{prooftree}
    \hypo{\TB{\Gamma, x_1 : A, x_2 : \TThunk{B}}{t_1}{B}}
    \hypo{\TB{\Gamma}{t_2}{B}}
    \hypo{\TB{\Gamma}{t_3}{\TList{A}}}
    \infer3{\TB{\Gamma}{\tfoldr{x_1}{x_2}{t_1}{t_2}{t_3}}{B}}
  \end{prooftree}
\end{center}

\subsection{Semantics}

\begin{center}
  \begin{prooftree}
    \hypo{\gamma(x) = v}
    \infer1{\EL{\gamma}{x}{v}}
  \end{prooftree}
\end{center}

\begin{center}
  \begin{prooftree}
    \hypo{\EL{\gamma}{t_1}{v_1}}
    \hypo{\EL{\gamma, x \mapsto v_1}{t_2}{v_2}}
    \infer2{\EL{\gamma}{\tlet{x}{t_1}{t_2}}{v_2}}
  \end{prooftree}
\end{center}

\begin{center}
  \begin{prooftree}
    \infer0{\EL{\gamma}{\ttrue}{\strue}}
  \end{prooftree}
\end{center}

\begin{center}
  \begin{prooftree}
    \infer0{\EL{\gamma}{\tfalse}{\sfalse}}
  \end{prooftree}
\end{center}

\begin{center}
  \begin{prooftree}
    \hypo{\EL{\gamma}{t_1}{v}}
    \hypo{\EL{\gamma}{t_2}{v}}
    \infer2{\EL{\gamma}{\teq{t_1}{t_2}}{\strue}}
  \end{prooftree}
\end{center}

\begin{center}
  \begin{prooftree}
    \hypo{\EL{\gamma}{t_1}{v_1}}
    \hypo{\EL{\gamma}{t_2}{v_2}}
    \hypo{v_1 \ne v_2}
    \infer3{\EL{\gamma}{\teq{t_1}{t_2}}{\sfalse}}
  \end{prooftree}
\end{center}

\begin{center}
  \begin{prooftree}
    \hypo{\EL{\gamma}{t_1}{v_1}}
    \hypo{\EL{\gamma}{t_2}{v_2}}
    \hypo{\sle{v_1}{v_2}}
    \infer3{\EL{\gamma}{\tle{t_1}{t_2}}{\strue}}
  \end{prooftree}
\end{center}

\begin{center}
  \begin{prooftree}
    \hypo{\EL{\gamma}{t_1}{v_1}}
    \hypo{\EL{\gamma}{t_2}{v_2}}
    \hypo{\snle{v_1}{v_2}}
    \infer3{\EL{\gamma}{\tle{t_1}{t_2}}{\sfalse}}
  \end{prooftree}
\end{center}

\begin{center}
  \begin{prooftree}
    \hypo{\EL{\gamma}{t_1}{\strue}}
    \hypo{\EL{\gamma}{t_2}{v}}
    \infer2{\EL{\gamma}{\tif{t_1}{t_2}{t_3}}{v}}
  \end{prooftree}
\end{center}

\begin{center}
  \begin{prooftree}
    \hypo{\EL{\gamma}{t_1}{\sfalse}}
    \hypo{\EL{\gamma}{t_3}{v}}
    \infer2{\EL{\gamma}{\tif{t_1}{t_2}{t_3}}{v}}
  \end{prooftree}
\end{center}

\begin{center}
  \begin{prooftree}
    \hypo{\EL{\gamma}{t_1}{v_1}}
    \hypo{\EL{\gamma}{t_2}{v_2}}
    \infer2{\EL{\gamma}{\tpair{t_1}{t_2}}{(v_1, v_2)}}
  \end{prooftree}
\end{center}

\begin{center}
  \begin{prooftree}
    \hypo{\EL{\gamma}{t}{v}}
    \infer1{\EL{\gamma}{\tthunk{t}}{\sthunk{v}}}
  \end{prooftree}
\end{center}

\begin{center}
  \begin{prooftree}
    \infer0{\EL{\gamma}{\tundefined}{\sundefined}}
  \end{prooftree}
\end{center}

\begin{center}
  \begin{prooftree}
    \hypo{\EL{\gamma}{t_1}{\sthunk{v_1}}}
    \hypo{\EL{\gamma, x \mapsto v_1}{t_2}{v_2}}
    \infer2{\EL{\gamma}{\tcase{t_1}{x}{t_2}{t_3}}{v_2}}
  \end{prooftree}
\end{center}

\begin{center}
  \begin{prooftree}
    \hypo{\EL{\gamma}{t_1}{\sundefined}}
    \hypo{\EL{\gamma}{t_3}{v}}
    \infer2{\EL{\gamma}{\tcase{t_1}{x}{t_2}{t_3}}{v}}
  \end{prooftree}
\end{center}

\begin{center}
  \begin{prooftree}
    \infer0{\EL{\gamma}{n}{n}}
  \end{prooftree}
\end{center}

\begin{center}
  \begin{prooftree}
    \hypo{\EL{\gamma}{t_1}{n_1}}
    \hypo{\EL{\gamma}{t_2}{n_2}}
    \infer2{\EL{\gamma}{\tadd{t_1}{t_2}}{n_1 + n_2}}
  \end{prooftree}
\end{center}

\begin{center}
  \begin{prooftree}
    \infer0{\EL{\gamma}{\tnil}{\snil}}
  \end{prooftree}
\end{center}

\begin{center}
  \begin{prooftree}
    \hypo{\EL{\gamma}{t_1}{v_1}}
    \hypo{\EL{\gamma}{t_2}{v_2}}
    \infer2{\EL{\gamma}{\tcons{t_1}{t_2}}{\scons{v_1}{v_2}}}
  \end{prooftree}
\end{center}

\begin{center}
  \begin{prooftree}
    \hypo{\EL{\gamma}{t_3}{v_1}}
    \hypo{\ELF{\gamma}{x_1}{x_2}{t_1}{t_2}{v_1}{v_2}}
    \infer2{\EL{\gamma}{\tfoldr{x_1}{x_2}{t_1}{t_2}{t_3}}{v_2}}
  \end{prooftree}
\end{center}

\begin{center}
  \begin{prooftree}
    \hypo{v : A}
    \infer1{\EL{\gamma}{\tfree{A}}{v}}
  \end{prooftree}
\end{center}

\begin{center}
  \begin{prooftree}
    \hypo{\EL{\gamma}{t_1}{v}}
    \infer1{\EL{\gamma}{\tchoose{t_1}{t_2}}{v}}
  \end{prooftree}
\end{center}

\begin{center}
  \begin{prooftree}
    \hypo{\EL{\gamma}{t_2}{v}}
    \infer1{\EL{\gamma}{\tchoose{t_1}{t_2}}{v}}
  \end{prooftree}
\end{center}

As in clairvoyance semantics, the semantics of \( \tfoldrZ \) is defined in terms of an auxiliary relation.

\begin{center}
  \begin{prooftree}
    \hypo{\EL{\gamma}{t_2}{v}}
    \infer1{\ELF{\gamma}{x_1}{x_2}{t_1}{t_2}{\snil}{v}}
  \end{prooftree}
\end{center}

\begin{center}
  \begin{prooftree}
    \hypo{\EL{\gamma, x_1 \mapsto v_1, x_2 \mapsto \sundefined}{t_1}{v}}
    \infer1{\ELF{\gamma}{x_1}{x_2}{t_1}{t_2}{\scons{v_1}{\sundefined}}{v}}
  \end{prooftree}
\end{center}

\begin{center}
  \begin{prooftree}
    \hypo{\ELF{\gamma}{x_1}{x_2}{t_1}{t_2}{\sthunk{v_2}}{v_3}}
    \hypo{\EL{\gamma, x_1 \mapsto v_1, x_2 \mapsto \sthunk{v_3}}{t_1}{v}}
    \infer2{\ELF{\gamma}{x_1}{x_2}{t_1}{t_2}{\scons{v_1}{v_2}}{v}}
  \end{prooftree}
\end{center}

\begin{lemma}
  If \( \TL{\Gamma}{t}{A} \), \( \gamma \in \Approx{\Gamma} \), and \( \EL{\gamma}{t}{v} \), then \( v \in \Approx{A} \).
\end{lemma}

Note that \( \tfail \) never reduces, so the language is not total.

\subsection{The clairvoyance translation}

We now define some constructs on the terms in \( \LL \).  For each construct \( c \), we provide typing rules and a set of evaluation rules that \emph{completely characterize} \( c \).  By ``completely characterize'', we mean the following.

\begin{itemize}
\item \emph{Soundness:} Each rule is valid.
\item \emph{Completeness:} Any derivation of the form \( \EL{\gamma}{c}{v} \) must have the form of one of the elimation rules for that construct.  In particular, there must be some elimination rule whose hypotheses are all true.  This allows us to ``invert'' judgments of the form \( \EL{\gamma}{c}{v} \).
\end{itemize}

\subsubsection{\texttt{return}}

\[
  \texttt{\treturn{t}} = \tpair{t}{0}
\]

\begin{center}
  \begin{prooftree}
    \hypo{\TL{\Gamma}{t}{A}}
    \infer1{\TL{\Gamma}{\treturn{t}}{\TM{A}}}
  \end{prooftree}
\end{center}

\begin{center}
  \begin{prooftree}
    \hypo{\EL{\gamma}{t}{v}}
    \infer1{\EL{\gamma}{\treturn{t}}{\spair{v}{0}}}
  \end{prooftree}
\end{center}

\subsubsection{bind}

\[
  \texttt{do $x$ <- $t_1$; $t_2$} = \tlet{\tpair{x}{y_1}}{t_1}{\tlet{\tpair{y_2}{y_3}}{t_2}{\tpair{y_2}{y_1 + y_3}}}
\]

\begin{center}
  \begin{prooftree}
    \hypo{\TL{\Gamma}{t_1}{\TM{A}}}
    \hypo{\TL{\Gamma, x : A}{t_2}{\TM{B}}}
    \infer2{\TL{\Gamma}{\texttt{do $x$ <- $t_1$; $t_2$}}{\TM{B}}}
  \end{prooftree}
\end{center}

\subsubsection{\texttt{transpose}}

\[
\begin{array}{l@{}l@{}l@{}}
  \texttt{transpose \(t\)} = \texttt{case \(t\) of } & \texttt{thunk \(y_1\) -> do $y_2$ <- $y_1$; return (\tthunk{y_2})} \\ & \texttt{undefined -> return undefined}
\end{array}
\]

\begin{center}
  \begin{prooftree}
    \hypo{\TL{\Gamma}{t}{\TThunk{(\TM{A})}}}
    \infer1{\TL{\Gamma}{\ttranspose{t}}{\TM{(\TThunk{A})}}}
  \end{prooftree}
\end{center}

\begin{center}
  \begin{prooftree}
    \hypo{\EL{\gamma}{t}{\sundefined}}
    \infer1{\EL{\gamma}{\ttranspose{t}}{\spair{\sundefined}{0}}}
  \end{prooftree}
\end{center}

\begin{center}
  \begin{prooftree}
    \hypo{\EL{\gamma}{t}{\sthunk{\spair{v}{c}}}}
    \infer1{\EL{\gamma}{\ttranspose{t}}{\spair{\sthunk{v}}{c}}}
  \end{prooftree}
\end{center}

\subsubsection{\texttt{foldM}}

\[
  \texttt{\tfoldM{x_1}{x_2}{t_1}{t_2}{t_3} = \tfoldr{x_1}{y_2}{\texttt{do \(x_2\) <- \ttranspose{y_2}; \(t_1\)}}{t_2}{t_3}} .
\]

\begin{center}
  \begin{prooftree}
    \hypo{\TL{\Gamma, x_1 : A, x_2 : \TThunk{B}}{t_1}{\TM{B}}}
    \hypo{\TL{\Gamma}{t_2}{\TM{B}}}
    \hypo{\TL{\Gamma}{t_3}{\TList{A}}}
    \infer3{\TL{\Gamma}{\texttt{foldM(\(\lambda x_1 x_2\).\(t_1\), \(t_2\), \(t_3\))}}{\TM{B}}}
  \end{prooftree}
\end{center}

To characterize the reduction behavior of \( \rulename{L-FoldM} \), we define yet another auxiliary relation:
\[
  \ELFM{\gamma}{x_1}{x_2}{t_1}{t_2}{v}{u} \equiv \ELF{\gamma}{x_1}{y_2}{\texttt{do \(x_2\) <- \ttranspose{y_2}; \(t_1\)}}{t_2}{v}{u} .
\]

\begin{lemma}
  The following rules completely characterize \( \tfoldMZ \).
  \begin{center}
    \begin{prooftree}
      \hypo{\ELFM{\gamma}{x_1}{x_2}{t_1}{t_2}{v}{u}}
      \infer1[\rulename{L-FoldM}]{\EL{\gamma}{\tfoldM{x_1}{x_2}{t_1}{t_2}{v}}{u}}
    \end{prooftree}
  \end{center}\begin{center}
    \begin{prooftree}
      \hypo{\EL{\gamma}{t_2}{u}}
      \infer1[\rulename{L-FoldM-Nil}]{\ELFM{\gamma}{x_1}{x_2}{t_1}{t_2}{\varepsilon}{u}}
    \end{prooftree}
  \end{center}
  \begin{center}
    \begin{prooftree}
      \hypo{\EL{\gamma, x_1 \mapsto v_1, x_2 \mapsto \sundefined}{t_1}{u}}
      \infer1[\rulename{L-FoldM-Undefined}]{\ELFM{\gamma}{x_1}{x_2}{t_1}{t_2}{\scons{v_1}{\sundefined}}{u}}
    \end{prooftree}
  \end{center}
  \begin{center}
    \begin{prooftree}
      \hypo{\ELFM{\gamma}{x_1}{x_2}{t_1}{t_2}{v_2}{\spair{v_3}{c_1}}}
      \hypo{\EL{\gamma, x_1 \mapsto v_1, x_2 \mapsto \sthunk{v_2}}{t_1}{\spair{v_4}{c_2}}}
      \infer2[\rulename{L-FoldM-Thunk}]{\ELFM{\gamma}{x_1}{x_2}{t_1}{t_2}{\scons{v_1}{\sthunk{v_2}}}{\spair{v_4}{c_1 + c_2}}}
    \end{prooftree}
  \end{center}
\end{lemma}
\begin{proof}
  We first prove validity.
  \newpage
  \newgeometry{margin=1cm}
  \begin{landscape}
    \begin{itemize}
    \item \textbf{\rulename{L-FoldM}.}
      \begin{center}
        \begin{prooftree}
          \hypo{\ELF{\gamma}{x_1}{y_2}{\texttt{do \(x_2\) <- \ttranspose{y_2}; \(t_1\)}}{t_2}{v}{u}}
          \infer1[\rulename{L-Foldr}]{\EL{\gamma}{\tfoldr{x_1}{y_2}{\texttt{do \(x_2\) <- \ttranspose{y_2}; \(t_1\)}}{t_2}{v}}{u}}
        \end{prooftree}
      \end{center}
    \item \textbf{\rulename{L-FoldM-Nil}.}
      \begin{center}
        \begin{prooftree}
          \hypo{\EL{\gamma}{t_2}{u}}
          \infer1[\rulename{L-Foldr-Nil}]{\ELF{\gamma}{x_1}{y_2}{\texttt{do \(x_2\) <- \ttranspose{y_2}; \(t_1\)}}{t_2}{\varepsilon}{u}}
        \end{prooftree}
      \end{center}
    \item \textbf{\rulename{L-FoldM-Undefined}.}
      \begin{center}
        \begin{prooftree}
          \infer0[\rulename{L-Var}]{\EL{\gamma, x_1 \mapsto v_1, y_2 \mapsto \bot}{y_2}{\bot}}
          \infer1[\rulename{L-Transpose-Undefined}]{\EL{\gamma, x_1 \mapsto v_1, y_2 \mapsto \bot}{\ttranspose{y_2}}{\spair{\bot}{0}}}
          \hypo{\EL{\gamma, x_1 \mapsto v_1, y_2 \mapsto \bot, x_2 \mapsto \sundefined}{t_1}{u}}
          \infer2[\rulename{L-Bind}]{\EL{\gamma, x_1 \mapsto v_1, y_2 \mapsto \bot}{\texttt{do \(x_2\) <- \ttranspose{y_2}; \(t_1\)}}{u}}
          \infer1[\rulename{L-Foldr-Undefined}]{\ELF{\gamma}{x_1}{y_2}{\texttt{do \(x_2\) <- \ttranspose{y_2}; \(t_1\)}}{t_2}{\scons{v_1}{\sundefined}}{u}}
        \end{prooftree}
      \end{center}
    \item \textbf{\rulename{L-FoldM-Thunk}.}  By induction, \( \ELF{\gamma}{x_1}{y_2}{\texttt{do \(x_2\) <- \ttranspose{y_2}; \(t_1\)}}{t_2}{v_2}{v_3} \).  Then,
      \begin{center}
        \begin{prooftree}
          \hypo{\ELF{\gamma}{x_1}{y_2}{\texttt{do \(x_2\) <- \ttranspose{y_2}; \(t_1\)}}{t_2}{v_2}{v_3}}
          \infer0[\rulename{L-Var}]{\EL{\gamma, x_1 \mapsto v_1, y_2 \mapsto \sthunk{\spair{v_3}{c_1}}}{y_2}{\sthunk{\spair{v_3}{c_1}}}}
          \infer1[\rulename{L-Transpose-Thunk}]{\EL{\gamma, x_1 \mapsto v_1, y_2 \mapsto \sthunk{\spair{v_3}{c_1}}}{\ttranspose{y_2}}{\spair{\sthunk{v_3}}{c_1}}}
          \hypo{\EL{\gamma, x_1 \mapsto v_1, y_2 \mapsto \sthunk{\spair{v_3}{c_2}}, x_2 \mapsto \sthunk{v_2}}{t_1}{\spair{v_4}{c_2}}}
          \infer2[\rulename{L-Bind}]{\EL{\gamma, x_1 \mapsto v_1, y_2 \mapsto \sthunk{\spair{v_3}{c_2}}}{\texttt{do \(x_2\) <- \ttranspose{y_2}; \(t_1\)}}{\spair{v_4}{c_1 + c_2}}}
          \infer2[\rulename{L-Foldr-Thunk}]{\ELF{\gamma}{x_1}{y_2}{\texttt{do \(x_2\) <- \ttranspose{y_2}; \(t_1\)}}{t_2}{\scons{v_1}{\sthunk{v_2}}}{\spair{v_4}{c_1 + c_2}}}
        \end{prooftree}
      \end{center}
    \end{itemize}
    To prove completeness, repeatedly apply inversion principles to find that any derivation of \( \tfoldM{x_1}{x_2}{t_1}{t_2}{v}{u} \) must have the form of one of the preceding derivations.
  \end{landscape}
  \restoregeometry
\end{proof}

\begin{center}
  \begin{align*}
    \floor{\tfoldr{x_1}{x_2}{t_1}{t_2}{t_3}} &= \texttt{do $y_3$ <- \(\floor{t_3}\); \tfoldM{x_1}{x_2}{\floor{t_1}}{\floor{t_2}}{y_3}}
  \end{align*}
\end{center}

\subsubsection{Properties of the translation}

\begin{lemma}
  If \( \TB{\Gamma}{t}{A} \), then \( \TL{\Gamma}{\floor{t}}{A} \).
\end{lemma}

\section{Equivalence}

\subsubsection{Clairvoyance semantics}

\begin{theorem}
  \( \EC{\gamma}{t}{c}{v} \) if and only if \( \EL{\gamma}{\floor{t}}{(v, c)} \).
\end{theorem}
\begin{proof}
  (\(\implies\)) By induction on the derivation of \( \EC{\gamma}{t}{c}{v} \).

  \textbf{Case:}
  \begin{center}
    \begin{prooftree}
      \hypo{\EC{\gamma}{t_3}{c_1}{v_1}}
      \hypo{\ECF{\gamma}{x_1}{x_2}{t_1}{t_2}{v_1}{c_2}{v_2}}
      \infer2[\rulename{L-Foldr}]{\EC{\gamma}{\tfoldr{x_1}{x_2}{t_1}{t_2}{t_3}}{c_1 + c_2}{v_2}}
    \end{prooftree}
  \end{center}
  \begin{itemize}
  \item \( t = \tfoldr{x_1}{x_2}{t_1}{t_2}{t_3} \)
  \item \( \floor{t} = \texttt{do \(y_3\) <- \(\floor{t_3}\); \tfoldM{x_1}{x_2}{\floor{t_1}}{\floor{t_2}}{y_3}} \)
  \end{itemize}
  We need to construct a derivation of the form
  \begin{center}
    \begin{prooftree}
      \hypo{\EL{\gamma}{\floor{t_3}}{\spair{v_1}{c_1}}}
      \hypo{\EL{\gamma, y_3 \mapsto v_1}{\tfoldM{x_1}{x_2}{\floor{t_1}}{\floor{t_2}}{y_3}}{\spair{v_2}{c_2}}}
      \infer2[\rulename{L-Bind}]{\EL{\gamma}{\texttt{do \(y_3\) <- \(\floor{t_3}\); \tfoldM{x_1}{x_2}{\floor{t_1}}{\floor{t_2}}{y_3}}{\spair{v_2}{c_2}}}{\spair{v_2}{c_1 + c_2}}}
    \end{prooftree}
  \end{center}
  By the inductive hypothesis, we have
  \[
    \EL{\gamma}{\floor{t_3}}{(v_1, c_1)} ,
  \]
  so we just need to prove
  \[
    \EL{\gamma, y_3 \mapsto v_1}{\tfoldM{x_1}{x_2}{\floor{t_1}}{\floor{t_2}}{y_3}}{\floor{t_2}}{y_3}{\spair{v_2}{c_2}} .
  \]
  For this, it suffices to prove
  \[
    \ELF{\gamma}{x_1}{y_2}{\tfoldM{x_1}{x_2}{\floor{t_1}}{\floor{t_2}}{y_3}}{\floor{t_2}}{y_3}{\spair{v_2}{c_2}} .
  \]
  We must also prove the following lemma by mutual induction.
  \begin{lemma}
    \( \ECF{\gamma}{x_1}{x_2}{t_1}{t_2}{v_1}{c}{v_2} \) if and only if \( \ELFM{\gamma}{x_1}{x_2}{\floor{t_1}}{\floor{t_2}}{v_1}{\spair{v_2}{c}} \).
  \end{lemma}
  \begin{proof}
    By induction on the derivation of \( \ECF{\gamma}{x_1}{x_2}{t_1}{t_2}{v_1}{c}{v_2} \).
  \begin{itemize}
  \item \textbf{Case \rulename{C-Foldr-Nil}:}
    \begin{itemize}
    \item \( v_1 = \varepsilon \)
    \item \( \EC{\gamma}{t_2}{c}{v_2} \)
    \end{itemize}
    By mutual induction, we obtain \( \EL{\gamma}{\floor{t_2}}{\spair{v_2}{c}} \), so we can construct
    \begin{center}
      \begin{prooftree}
        \hypo{\EL{\gamma}{\floor{t_2}}{\spair{v_2}{c}}}
        \infer1[\rulename{L-FoldM-Nil}]{\ELF{\gamma}{x_1}{x_2}{\floor{t_1}}{\floor{t_2}}{\snil}{\spair{v_2}{c}}}
      \end{prooftree}
    \end{center}
  \item \textbf{Case \rulename{L-Foldr-Undefined}}:
  \end{itemize}

  (\(\impliedby\)) By induction on the derivation of \( \ELFM{\gamma}{x_1}{x_2}{\floor{t_1}}{\floor{t_2}}{v_1}{\spair{v_2}{c}} \).
  % \end{proof}
  \end{proof}
\end{proof}

\subsubsection{Demand semantics}

\begin{align*}
  \transD{t}{x_1 : A_1, \ldots, x_n : A_n} = &\texttt{let \( x^A_1 \) = \tfree{A_1} in} \\
                                             &\texttt{assert \tle{x^A_1}{x_1} in} \\
                                             &\vdots \\
                                             &\texttt{let \( x^A_n \) = \tfree{A_n} in} \\
                                             &\texttt{assert \tle{x^A_n}{x_n} in} \\
                                             &\texttt{let \( c \) = \tfree{\TNat} in} \\
                                             &\texttt{assert \teq{\transC{t}}{\tpair{y^A}{c}} in} \\
                                             &\texttt{(\( x^A_1 \), \( \ldots \), \( x^A_n \), \( c \))}
\end{align*}

\begin{lemma}
  Let \( \Gamma = x_1 : A_1, \ldots, x_n : A_n \).  If \( \TB{\Gamma}{t}{B} \) and \( \gamma \in \Eval{\Gamma} \), then \( \TL{\Gamma, y^A : B}{\transD{t}{\Gamma}}{\TM{(\texttt{\(A_1\)\ *\ \(\cdots\) *\ \(A_n\)})}} \)
\end{lemma}

\begin{theorem}
  If \( \ED{t}{\gamma}{b^A} = (\gamma^A, c) \), then \( \EL{\gamma, y^A \mapsto b^A}{\transD{t}{\Gamma}}{\spair{\verts{\gamma^A}}{c}} \), where \( \verts{x_1 \mapsto a_1, \ldots, x_n \mapsto a_n} = (a_1, \ldots, a_n) \).
\end{theorem}

\end{document}
